% =====================================================================
%  PR/FAQ - Diagonality: The Best of Both Worlds
%  AWS World Sports Innovation Cup 2026 - Challenge 2
%  Compile with Tectonic:  tectonic prfaq.tex
% =====================================================================
\documentclass[11pt]{article}

\usepackage[utf8]{inputenc}
\usepackage[T1]{fontenc}
\usepackage{lmodern}
\usepackage[a4paper,margin=2.5cm]{geometry}
\usepackage{graphicx}
\usepackage{xcolor}
\usepackage{titlesec}
\usepackage{enumitem}
\usepackage{parskip}
\usepackage[colorlinks=true,linkcolor=black,urlcolor=blue]{hyperref}

\definecolor{accent}{HTML}{1F6FB2}
\graphicspath{{figures/}}

\titleformat{\section}{\large\bfseries\color{accent}}{}{0pt}{}
\titleformat{\subsection}{\normalsize\bfseries}{}{0pt}{}
\titlespacing{\section}{0pt}{14pt}{6pt}
\titlespacing{\subsection}{0pt}{10pt}{3pt}

\newcommand{\faq}[1]{\subsection{#1}}

\begin{document}

% ---------------------------------------------------------------------
\begin{center}
{\LARGE\bfseries Diagonality: The Best of Both Worlds}\\[6pt]
{\large\color{accent} PR\,/\,FAQ}\\[4pt]
{\small AWS World Sports Innovation Cup 2026, Challenge 2: Unlock the
Power of 3D Football Data}\\[2pt]
{\small Jaime Oriol \quad$\cdot$\quad
\href{https://github.com/jaime-oriol/Diagonality_3D}{github.com/jaime-oriol/Diagonality\_3D}}
\end{center}

\hrule
\vspace{10pt}

% =====================================================================
\section{Press Release}

\textbf{\large Football's most-discussed tactical idea finally has a
number, and a real-time map.}

\textit{Diagonality, the preference for oblique passes, carries and runs
over straight ones, has been argued about for decades but never measured.
A new framework built on 3D skeleton tracking validates the theory against
6{,}923 on-ball actions and turns it into the Diagonal Opportunity Surface:
a live, orientation-aware map of where, when and against which defender to
attack diagonally.}

\vspace{6pt}
\textbf{Frankfurt, 16 May 2026.} Coaches have a name for the smartest
ball a team can play, the diagonal, and a theory for why it works:
Spielverlagerung's 2025 essay calls it ``the best of both worlds'', a pass
that carries the progression of a vertical ball and the safety of a
horizontal one. The essay also hedges every claim it makes, bracketing the
word ``(assumed)'', because diagonality is at bottom a statement about
\emph{body orientation}, and tracking data has only ever recorded
\emph{positions}. The idea could be admired but not tested, and certainly
not turned into a tool.

\textbf{Diagonality} closes that gap. Using TRACAB 3D skeleton tracking
from five Bundesliga 2025--26 matches (21 keypoints per player at
50\,Hz, the first dataset that measures the real head, shoulder and hip
orientation of every footballer), the project does two things. First, it
\textbf{tests the theory}: the essay's argument is distilled into a
perceptual premise and three falsifiable hypotheses, tested against
6{,}923 passes, carries and take-ons with metrics that owe nothing to the
new model. The claims hold.
Diagonal actions are the efficient compromise between progression and
safety, they disrupt the defensive block on the lateral axis a vertical
pass leaves untouched, and they hand the receiver a body already turned
toward goal.

Second, it \textbf{turns the validated theory into a tool}: the
\textbf{Diagonal Opportunity Surface (DOS)}. DOS is a real-time pitch map
that fuses what every defender can see, how each is oriented, and who
controls space, into a single signal: the extra control an attacker buys
by delivering the ball diagonally rather than straight, into the space
behind a defender who cannot see it. A high-DOS zone is not a black box: it
decomposes into the three causes the tactical theory names: a defender's
blind spot, the orientation delay in reaching it, and the directional
benefit of arriving oblique.

``Diagonality was the game's most-discussed unquantified idea, and it
stayed unquantified for a concrete reason: it is a claim about body
orientation, and tracking never recorded orientation,'' said Jaime Oriol,
who built the framework. ``With 3D skeleton data we could finally measure
it. We tested the theory and it held; then we turned it into a map a coach
can read frame by frame. We turned a sentence into a number on a pitch.''

For a coaching or analysis desk, DOS is immediately usable. Aggregated per
opponent, it is a scouting map of where a specific team's defenders
habitually turn away from danger. Aggregated per player, it is a profile,
and it passes the analysts' informal ``Messi test'': ranked by the total
diagonal opportunity they generate, the leaderboard is headed by exactly
the players the eye already knows as creators. Gated to what the player on
the ball can actually perceive, DOS becomes a real-time decision aid,
showing which passes, dribbles and off-ball runs are both available and
visible, plus a second amber layer of high-value options the carrier has
stopped watching and should turn to re-check.

``The honest problem with `play more diagonals' as an instruction is that
nobody could point to where or when,'' said a hypothetical performance
analyst, the kind of user the framework is built for. ``A surface that
lights up the blind side, per opponent, turns a slogan into something you
can coach on Monday and check on Saturday.''

The framework is released as open source. Every stage runs from the
repository root, is memory-safe and deterministic, and is documented for
reproduction; no proprietary hackathon data is redistributed. The same
pipeline scales without modification from five matches to a full season.

% =====================================================================
\section{Frequently Asked Questions: External}

\faq{What exactly is the Diagonal Opportunity Surface?}
DOS is a value, computed for every cell of the pitch at every frame, equal
to the extra ball control the attacking team would gain by playing into
that cell \emph{diagonally} rather than straight. It is built from a
per-defender vision model, an orientation-aware pitch-control function and
a cognitive scanning gate. The same routine scores a pass, a carry, a
take-on or an off-ball run, because all four reduce to an origin and a
direction.

\faq{How is this different from existing pitch-control models?}
Every pitch-control model in the literature treats a player as a disc or an
isotropic Gaussian, equally fast and equally aware in all directions.
None encodes which way the player is facing. DOS is built on the opposite
premise: a defender controls the space in front of them far better than the
space behind their shoulder, because their \emph{perception} of space is
asymmetric. That asymmetry is the entire mechanism of diagonality, and it
requires measured orientation to model.

\faq{Why 3D skeleton data and not ordinary tracking?}
Ordinary tracking gives positions; body orientation then has to be
proxied from the velocity vector, which a player violates at will by
running one way while facing another, an error of $30$--$50^\circ$, large
enough to drown the signal. Broadcast pose estimation reduces it to about
$27^\circ$. TRACAB 3D skeleton tracking removes it: the head, shoulder and
hip keypoints are measured directly, at 50\,Hz.

\faq{Did the tactical theory actually hold up?}
Yes, on three independent tests. Diagonal actions produce a value-positive
outcome more often than forward or sideways actions (the
progression--safety trade-off); they carry the highest \emph{combined}
longitudinal-plus-lateral disruption of the defensive block; and their
receiver needs far less rotation to face goal and sees more team-mates at
the moment of reception. Honest nuance: the diagonal does not beat the
vertical pass on \emph{depth} disruption: forward actions match it there.
The diagonal's specific edge is the lateral axis, and the combination.

\faq{Is DOS biased by hindsight, or does it secretly use outcomes?}
No, and this is deliberate. DOS is computed from geometry, the vision model
and skeleton orientation alone. It never sees a goal, an xG, an xT or a
success label. That independence is what makes the validation fair: when
DOS predicts expected-threat gain, it is predicting an outcome it was never
shown.

\faq{Can it run live during a match?}
The model is single-frame and the scanning gate is explicitly real-time:
it bounds the surface to what the on-ball player can perceive within a
2.5-second memory window. The current implementation favours clarity and
reproducibility over latency optimisation; a production deployment would
target the broadcast or bench-tablet frame budget.

\faq{Does it work for off-ball runs?}
Yes. Run detection is already a solved, commercial product; the framework
uses a lightweight detector and applies the same DOS valuation. 2{,}354
off-ball runs were detected and scored, and the overwhelming majority carry
positive DOS: off-ball movement, like on-ball action, is a search for the
defender's blind side.

\faq{What can a club do with it tomorrow?}
Four things: match preparation (a per-opponent map of orientation
vulnerabilities), post-match analysis (every action annotated with its DOS
and the misalignment it exploited), recruitment (players profiled by the
diagonal opportunity they generate), and broadcast (the blind-spot and DOS
surfaces as legible overlays).

% =====================================================================
\section{Frequently Asked Questions: Internal}

\faq{Why a reach-field pitch-control model instead of standard
time-to-intercept?}
Because the question is not ``who gets there first'' but ``how well is each
player placed, given how they are turned''. Each player is modelled as an
anisotropic Gaussian whose width follows an orientation-aware delay: a
reaction-time term that grows with visual eccentricity plus a
change-of-direction penalty, both taken from published biomechanics. A
defender with a threat behind their shoulder has a control blob compressed
to half-width: a literal hole in the place the theory says the hole should
be.

\faq{What was the hardest part?}
Making a perceptual idea computable without it becoming arbitrary. The
vision model, the awareness weighting and the detection delay each had to
trace back to a published mechanism rather than a tuned constant, so that a
high-DOS cell could be \emph{explained}, not just displayed.

\faq{What failed along the way, and what did it teach you?}
Several things, and each shaped the final design:
\begin{itemize}[leftmargin=1.3em,topsep=2pt,itemsep=2pt]
  \item \textbf{Memory.} A single match does not fit in 8\,GB of RAM. The
  naive ``load the match'' approach failed immediately; every stage was
  rebuilt around chunked, predicate-pushdown reads over frame ranges. The
  pipeline is now memory-safe by construction and, as a side effect,
  trivially scalable.
  \item \textbf{Carries.} Scoring a carry by its start-to-end direction was
  wrong: a carry serpentines, stops and turns. The fix was to sample the
  carry at several frames and use the \emph{instantaneous} skeleton
  velocity as the true direction, taking the diagonal peak.
  \item \textbf{A data quirk.} The DFL event XML does not list carries
  under their parent possession, which silently broke outcome linkage. It
  had to be repaired by matching on frame range and team.
  \item \textbf{Visual flicker.} The first real-time DOS render flickered
  frame to frame. The cause was relative thresholds; the fix was absolute
  thresholds plus a temporal smoothing pass with a hard reset on possession
  change.
  \item \textbf{A surprising result.} DOS and the defensive-disruption
  metric D-Def turned out only weakly (and negatively) correlated. The
  first instinct was that something was broken. It was not: DOS measures the
  perceptual \emph{opportunity} at the instant of the action, D-Def the
  structural \emph{consequence} three seconds later. They are complementary,
  and the gap between them is itself a coaching signal.
\end{itemize}

\faq{Why the ``Messi test'', and did it pass?}
A metric that claims to capture a footballing quality must rank the players
the eye already knows for it, or it is measuring something else. Ranked by
total diagonal opportunity generated, the leaderboard is headed by a
recognisable creative core; by per-action average, a centre-back and a
goalkeeper sit at the foot of the table. A surface built only from skeleton
geometry, with no notion of reputation, recovers the game's own hierarchy
of creators. It passed.

\faq{What are the real limitations?}
The sample is five matches, the data released for the challenge. This is
a limit on the \emph{input}, not the method: the pipeline is match-agnostic
and deterministic, so a full season is a matter of feeding in more data,
not changing code. The biomechanical constants are taken from the
literature rather than fitted to football tracking; fitting them on a
larger sample would sharpen the surface. And the orientation-aware
pitch-control function is currently validated indirectly, through DOS.

\faq{What did it cost to run?}
The metric chain and statistics run on a local 8\,GB machine thanks to the
memory-safe design. The video renders, being embarrassingly parallel, were
produced on AWS compute. No specialised hardware and no GPU are required
for the analysis itself.

\faq{What is next?}
A learned run-detection front-end so DOS scores every run, not only
sustained bursts; a direct comparison of the reach-field pitch-control
function against a standard model; biomechanical constants fitted on a
larger sample; and a latency-optimised build for live deployment.

\vspace{12pt}
\hrule
\vspace{6pt}
{\small The full technical write-up, all source code and reproduction
instructions are in the repository:
\href{https://github.com/jaime-oriol/Diagonality_3D}{github.com/jaime-oriol/Diagonality\_3D}.}

\end{document}
